% IEEE Conference Style Report
\documentclass[conference]{IEEEtran}

\usepackage{graphicx}
\usepackage{booktabs}
\usepackage{url}

\title{Local AI Dashcam Incident Explainer: An End-to-End, Privacy-Preserving Incident Analysis Pipeline}

\author{
    \IEEEauthorblockN{Kunal Tailor}
    \IEEEauthorblockA{Computer Vision Course Project\\
    \textit{Local AI Dashcam Incident Explainer}\\
    March 2026}
}

\begin{document}
\maketitle

\begin{abstract}
Dashcam footage is a valuable source for incident analysis, but manual review is slow and privacy sensitive. This report presents an end-to-end local pipeline that detects incident windows, samples keyframes, produces structured visual-language analysis, and generates a human-readable insurance report with a downloadable PDF. The system combines classical computer vision, object detection and tracking, optical flow, local vision-language models, and privacy filters. The result is a fully offline workflow that transforms raw driving video into a structured evidence package for review.
\end{abstract}

\begin{IEEEkeywords}
dashcam analytics, incident detection, optical flow, object tracking, vision-language models, privacy, local AI
\end{IEEEkeywords}

\section{Introduction}
Traffic incident review typically requires manual inspection of long, privacy-sensitive dashcam videos. We propose a fully local pipeline that compresses this workflow into a repeatable, auditable process: detect the incident window, summarize the scene with a vision-language model, and export a structured PDF report. The system is designed for coursework and real-world review workflows that require privacy preservation and offline execution.

\section{Related Work}
Our detection stack builds on real-time multi-object tracking and optical flow. We adopt a SORT-style tracking approach based on Kalman filtering and IoU-based assignment~\cite{bewley2016sort} and use Farneb\"ack dense optical flow for motion anomalies~\cite{farneback2003flow}. For vision-language narration, we use MiniCPM-V, an efficient multimodal model designed for local deployment~\cite{minicpmv2024}. Related benchmarks in traffic anomaly detection and accident anticipation include DoTA~\cite{dota2023}, DADA-2000~\cite{dada2019}, and uncertainty-based accident anticipation approaches~\cite{uncertainty2020}. Object detection is performed with Ultralytics YOLO models~\cite{ultralytics}.

\section{System Overview (End-to-End)}
The pipeline operates in five phases:
\begin{enumerate}
    \item \textbf{Keyframe Extraction}: Uniform or motion-triggered extraction reduces redundant frames.
    \item \textbf{Incident Detection}: YOLO detections are tracked with SORT; optical flow spikes signal anomalies.
    \item \textbf{VLM Narration}: Three keyframes (pre, impact, post) are summarized into structured JSON.
    \item \textbf{Privacy Filtering}: Faces and license plates are blurred before display or export.
    \item \textbf{Report Export}: JSON and keyframes are compiled into a PDF report and downloadable assets.
\end{enumerate}
The system is accessible via a local Gradio UI and a CLI pipeline for batch processing.

\section{Methodology}
\subsection{Keyframe Sampling}
Frames are extracted uniformly via FFmpeg or using a motion-triggered threshold based on frame differencing. This reduces compute while retaining the most informative context around an event.

\subsection{Incident Detection and Tracking}
Each sampled frame is passed through YOLO detection. A SORT-based tracker assigns stable IDs across frames, while Farneback optical flow measures motion magnitude. An incident is flagged when a high-IoU interaction coincides with a motion spike in the same temporal window.

\subsection{Vision-Language Narration}
Three representative frames (pre, impact, post) are encoded and passed to a local VLM. The model is prompted to return structured JSON fields including vehicles, severity, fault analysis, and timeline notes. A local LLM then generates the narrative report.

\subsection{Privacy and Export}
Faces are blurred using a DNN face detector, and license plates are blurred using vehicle ROI heuristics. The final evidence package includes the incident clip, annotated keyframes, JSON, and a PDF report.

\section{Implementation}
The system is implemented in Python with OpenCV, Ultralytics YOLO, and Gradio. Models are executed locally via Ollama. Outputs are written to timestamped folders under \texttt{outputs/} for auditability. The UI provides video upload, real-time logs, a keyframe gallery, JSON viewer, report text, and PDF download.

\section{Evaluation}
We evaluate report quality using BLEU-4 and ROUGE-L and study the tradeoff between keyframe count and inference time. An internal ablation compares 1, 3, and 5 keyframe settings. Three keyframes provide the best balance between speed and narrative completeness.

\section{Results}
Table~\ref{tab:ablation} summarizes the keyframe ablation trend observed in local testing.

\begin{table}[ht]
\caption{Keyframe Ablation (Local Validation)}
\label{tab:ablation}
\centering
\begin{tabular}{lccc}
\toprule
\textbf{Keyframes} & \textbf{Inference Time} & \textbf{ROUGE-L} & \textbf{Notes} \\
\midrule
1 & \textasciitilde18s & 0.28 & Misses pre/post context \\
3 & \textasciitilde42s & 0.41 & Best balance \\
5 & \textasciitilde71s & 0.43 & Marginal gain \\
\bottomrule
\end{tabular}
\end{table}

\section{Limitations}
The VLM can hallucinate in low-light scenes, and optical-flow spikes can be triggered by sudden camera shake. License plate blurring is heuristic and may miss extreme angles. Future work includes video-native VLMs, specialized plate detectors, and GPU-accelerated pipelines for reduced latency.

\section{Conclusion}
This project demonstrates a complete end-to-end pipeline that turns raw dashcam footage into a privacy-safe evidence package. By combining classic CV, detection and tracking, VLM narration, and report generation, the system supports offline, auditable incident analysis workflows suitable for coursework and operational review.

\begin{thebibliography}{00}
\bibitem{bewley2016sort} A. Bewley, Z. Ge, L. Ott, F. Ramos, and B. Upcroft, ``Simple Online and Realtime Tracking,'' arXiv:1602.00763, 2016. [Online]. Available: \url{https://arxiv.org/abs/1602.00763}
\bibitem{farneback2003flow} G. Farneb\"ack, ``Two-Frame Motion Estimation Based on Polynomial Expansion,'' in \textit{SCIA 2003}, LNCS 2749, pp. 363--370, 2003. [Online]. Available: \url{https://www.ida.liu.se/ext/WITAS-ev/Computer_Vision_Technologies/PaperInfo/farneback03.html}
\bibitem{minicpmv2024} Y. Yao \textit{et al.}, ``MiniCPM-V: A GPT-4V Level MLLM on Your Phone,'' arXiv:2408.01800, 2024. [Online]. Available: \url{https://arxiv.org/abs/2408.01800}
\bibitem{dota2023} Y. Yao \textit{et al.}, ``DoTA: Unsupervised Detection of Traffic Anomaly in Driving Videos,'' \textit{IEEE Trans. Pattern Anal. Mach. Intell.}, vol. 45, no. 1, pp. 444--459, 2023, doi: 10.1109/TPAMI.2022.3150763. [Online]. Available: \url{https://pubmed.ncbi.nlm.nih.gov/35157576/}
\bibitem{dada2019} J. Fang \textit{et al.}, ``DADA-2000: Can Driving Accident be Predicted by Driver Attention?'' arXiv:1904.12634, 2019. [Online]. Available: \url{https://arxiv.org/abs/1904.12634}
\bibitem{uncertainty2020} W. Bao, Q. Yu, and Y. Kong, ``Uncertainty-based Traffic Accident Anticipation with Spatio-Temporal Relational Learning,'' in \textit{Proc. ACM MM}, pp. 2682--2690, 2020, doi: 10.1145/3394171.3413827. [Online]. Available: \url{https://cir.nii.ac.jp/crid/1360301475480193664}
\bibitem{ultralytics} Ultralytics, ``Ultralytics YOLO Documentation.'' [Online]. Available: \url{https://docs.ultralytics.com/}
\end{thebibliography}

\end{document}
